\documentclass[a4paper,12pt]{article}
\usepackage[T2A]{fontenc}
\usepackage[utf8]{inputenc}
\usepackage[russian]{babel}
\usepackage{amsmath,amssymb,mathtools}
\usepackage{icomma}
\usepackage[a4paper,left=22mm,right=21mm,top=24mm,bottom=24mm,headheight=16pt,headsep=10pt]{geometry}
\usepackage{microtype}
\usepackage{xcolor}
\usepackage{tikz}
\usetikzlibrary{arrows.meta,positioning,shapes.geometric,calc}
\usepackage{booktabs,tabularx}
\usepackage{enumitem}
\usepackage{fancyhdr}
\usepackage{titlesec}
\usepackage{setspace}
\usepackage[hidelinks]{hyperref}
\definecolor{accent}{HTML}{007F80}
\definecolor{darkaccent}{HTML}{00595C}
\definecolor{textgray}{HTML}{343A3B}
\definecolor{linegray}{HTML}{BCC8C9}
\definecolor{lightaccent}{HTML}{EEF6F4}
\color{textgray}
\onehalfspacing
\setlength{\parindent}{0pt}
\setlength{\parskip}{5pt}
\setcounter{secnumdepth}{1}
\titleformat{\section}{\large\color{darkaccent}\normalfont}{\thesection.}{0.55em}{}
\titlespacing*{\section}{0pt}{15pt}{6pt}
\titleformat{\subsection}{\normalsize\color{darkaccent}\normalfont}{\thesubsection.}{0.55em}{}
\titlespacing*{\subsection}{0pt}{11pt}{4pt}
\setlist[enumerate]{leftmargin=*,itemsep=4pt,topsep=5pt}
\setlist[itemize]{leftmargin=*,itemsep=2pt,topsep=3pt}
\pagestyle{fancy}
\fancyhf{}
\fancyhead[L]{\small Отношения и пропорции}
\fancyhead[R]{\small Надя Клименко}
\fancyfoot[C]{\small\thepage}
\renewcommand{\headrulewidth}{0.35pt}
\renewcommand{\footrulewidth}{0pt}
\newcommand{\keyidea}[1]{\textcolor{darkaccent}{\textbf{Главная мысль.}} #1}
\newcommand{\example}[1]{\par\medskip\textcolor{darkaccent}{\textbf{Пример.}} #1\par}
\newcommand{\checkmarkline}[1]{\par\smallskip\textcolor{darkaccent}{$\square$}\ #1\par}
\hypersetup{pdftitle={Отношения и пропорции — чек-лист},pdfauthor={Лёвин Артём Александрович},pdfsubject={Математика; Надя Клименко; 21 ноября 2026}}
\begin{document}
\pagenumbering{gobble}

% Титульный лист
\thispagestyle{empty}
\begin{center}
\vspace*{23mm}
{\color{darkaccent}\Large МАТЕМАТИКА\par}
\vspace{24mm}
{\LARGE Чек-лист\par}
\vspace{14mm}
{\color{accent}\LARGE Отношения и пропорции\par}
\vspace{13mm}
{\large Часть и целое. Деление в отношении.\\
Основное свойство пропорции. Выражение переменной.\par}
\vfill
{\large Надя Клименко\par}
\vspace{14mm}
{21 ноября 2026 г.\par}
\vspace{24mm}
{\small Лёвин Артём Александрович\\
эксклюзивно для Нади Клименко\par}
\vspace{18mm}
\begin{tikzpicture}[x=1cm,y=1cm]
\draw[accent,line width=.75pt] (0,0) .. controls (1.5,.65) and (2.1,-.65) .. (3.6,0)
  .. controls (5.0,.65) and (5.8,-.65) .. (7.3,0);
\fill[darkaccent] (3.6,0) circle (1.05pt);
\end{tikzpicture}
\end{center}
\clearpage
\pagenumbering{arabic}

\section{Отношение двух величин}
Отношением числа $a$ к числу $b$ называют частное $a:b=\dfrac ab$, где $b\ne0$.
Порядок чисел важен: $a:b$ и $b:a$ обычно дают разные ответы.

\keyidea{Если требуется узнать, какую часть составляет одна величина от другой,
делите \emph{часть на целое}. Если требуется узнать, во сколько раз первая величина
больше второй, делите \emph{первую на вторую}.}

\example{Из пяти заданий выполнено четыре. Какую часть составляют выполненные?}
\[
\frac{4}{5}=0{,}8=80\%.
\]
Это три записи одного результата: обыкновенная дробь, десятичная дробь и проценты.

\example{Из $480\ \text{м}^2$ катка расчистили $320\ \text{м}^2$.}
\begin{align*}
\text{Расчищенная доля:}\quad &\frac{320}{480}=\frac23;\\
\text{Во сколько раз весь каток больше:}\quad &\frac{480}{320}=\frac32=1{,}5.
\end{align*}
Обратите внимание: в двух вопросах используются те же числа, но порядок деления меняется.

\medskip
\begin{tabularx}{\linewidth}{@{}X X@{}}
\toprule
\textcolor{darkaccent}{Формулировка} & \textcolor{darkaccent}{Что вычисляем}\\
\midrule
«Какую часть от целого составляет $A$?» & $A:\text{целое}$\\
«Во сколько раз $A$ больше $B$?» & $A:B$\\
«Отношение $A$ к $B$» & $A:B$\\
\bottomrule
\end{tabularx}

\checkmarkline{Могу определить, какую величину поставить в числитель и какую --- в знаменатель.}
\checkmarkline{Умею представить долю дробью, десятичным числом и процентами.}

\section{Одинаковые единицы измерения}
Перед сравнением физических величин приведите их к \textbf{одной единице измерения}.
Численное отношение однородных величин при этом не зависит от выбора единицы.

\example{Масса котёнка $99$ г, масса кошки $4{,}4$ кг. Найдите отношение массы котёнка к массе кошки.}
\[
4{,}4\ \text{кг}=4400\ \text{г},\qquad
\frac{99}{4400}=\boxed{\frac9{400}}=0{,}0225.
\]
Проверка разумности: масса котёнка меньше массы кошки, значит, отношение меньше единицы.
При переводе килограммов в граммы число умножают на $1000$.

\checkmarkline{Перед делением величин проверяю и согласую единицы измерения.}

\section{Деление числа в заданном отношении}
Запись $3:4$ означает, что первой величине соответствуют \textbf{три равные доли},
а второй --- \textbf{четыре такие доли}. Всего $3+4=7$ долей.

\example{В лагере $105$ детей, число девочек относится к числу мальчиков как $3:4$.}
\begin{align*}
&3+4=7 &&\text{долей всего},\\
&105:7=15 &&\text{детей в одной доле},\\
&3\cdot15=\boxed{45} &&\text{девочек},\\
&4\cdot15=\boxed{60} &&\text{мальчиков}.
\end{align*}
\begin{center}
\begin{tikzpicture}[x=1.18cm,y=1cm]
\foreach \i in {0,1,2} {\fill[lightaccent] (\i,0) rectangle ++(1,.62);}
\foreach \i in {3,4,5,6} {\fill[gray!7] (\i,0) rectangle ++(1,.62);}
\foreach \i in {0,...,6} {\draw[darkaccent,line width=.6pt] (\i,0) rectangle ++(1,.62);}
\node[font=\small,below] at (1.5,-.1) {девочки: $3\cdot15=45$};
\node[font=\small,below] at (5,-.1) {мальчики: $4\cdot15=60$};
\end{tikzpicture}
\end{center}
Проверка: $45+60=105$, а $45:60=3:4$.

\textcolor{darkaccent}{\textbf{Алгоритм.}} Чтобы разделить $S$ в отношении $p:q$:
\begin{enumerate}
\item Найдите число долей: $p+q$.
\item Найдите одну долю: $S:(p+q)$.
\item Умножьте величину доли сначала на $p$, затем на $q$.
\item Проверьте сумму и полученное отношение.
\end{enumerate}
Для $p:q:r$ действуйте так же, используя сумму $p+q+r$.

\example{Разделите $12$ в отношении $1:2:4$.}
\[
1+2+4=7,\quad \boxed{\frac{12}{7}},\quad
\boxed{\frac{24}{7}},\quad\boxed{\frac{48}{7}}.
\]
Доли могут быть дробными. Их сумма равна $12$.

\checkmarkline{Умею делить число на две или три части в указанном отношении.}

\section{Что такое пропорция}
\textbf{Пропорция} --- равенство двух отношений:
\[
\frac ab=\frac cd \qquad (b\ne0,\ d\ne0).
\]
То же можно записать через двоеточия: $a:b=c:d$.
В этой записи $a$ и $d$ называют \textbf{крайними членами},
а $b$ и $c$ --- \textbf{средними членами}.

\begin{center}
\begin{tikzpicture}[>=Latex]
\node (expr) at (0,0) {$\displaystyle a:b\;=\;c:d$};
\node[font=\small,text=darkaccent] at (-3.15,0.68) {крайние};
\node[font=\small,text=darkaccent] at (3.15,0.68) {средние};
\draw[->,darkaccent] (-2.25,.53) -- (-1.13,.2);
\draw[->,darkaccent] (-2.25,.53) -- (1.7,.18);
\draw[->,darkaccent] (2.2,.53) -- (-.5,.16);
\draw[->,darkaccent] (2.2,.53) -- (.88,.18);
\end{tikzpicture}
\end{center}
\keyidea{Основное свойство пропорции: произведение крайних членов
равно произведению средних.}
\[
\boxed{\frac ab=\frac cd\quad\Longleftrightarrow\quad ad=bc},
\qquad b\ne0,\ d\ne0.
\]

\example{Проверьте, верно ли равенство $\dfrac23=\dfrac46$.}
\[
2\cdot6=12,\qquad3\cdot4=12.
\]
Произведения равны, значит, пропорция верна.

\example{А равенство $\dfrac23=\dfrac45$?}
\[
2\cdot5=10,\qquad3\cdot4=12,\qquad10\ne12.
\]
Пропорция неверна.

\checkmarkline{Называю крайние и средние члены пропорции без подсказки.}
\checkmarkline{Проверяю пропорцию перемножением крайних и средних членов.}

\section{Нахождение неизвестного члена}
Если одна из четырёх величин неизвестна, примените основное свойство пропорции,
а затем решите получившееся уравнение.

\example{Найдите $t$ из пропорции $\dfrac{t}{4{,}2}=\dfrac{2{,}5}{3{,}5}$.}
\[
3{,}5t=4{,}2\cdot2{,}5=10{,}5,\qquad
\boxed{t=3}.
\]
Проверка: $3:4{,}2=2{,}5:3{,}5$.

\example{В одном из членов пропорции есть множитель:}
\[
\frac{13{,}7}{4}=\frac{9t}{3{,}6}.
\]
При перемножении считайте $9t$ одним выражением:
\[
13{,}7\cdot3{,}6=4\cdot9t,\qquad
49{,}32=36t,\qquad \boxed{t=1{,}37}.
\]

\example{В одном из членов пропорции есть сумма:}
\[
\frac{1{,}5}{6{,}3}=\frac{a+0{,}03}{0{,}21}.
\]
\begin{align*}
1{,}5\cdot0{,}21&=6{,}3(a+0{,}03),\\
0{,}315&=6{,}3(a+0{,}03),\\
a+0{,}03&=0{,}315:6{,}3=0{,}05,\\
a&=\boxed{0{,}02}.
\end{align*}
\textbf{Важно:} составное выражение при переходе к записи через двоеточие
сохраняют в скобках: $1{,}5:6{,}3=(a+0{,}03):0{,}21$.
Если сумма находится в знаменателе, например $\frac{a}{a+0{,}03}$,
получается $a:(a+0{,}03)$.

\checkmarkline{При перемножении сохраняю скобки вокруг составных выражений.}
\checkmarkline{После нахождения неизвестной проверяю исходную пропорцию.}

% Отдельная страница алгоритма
\clearpage
\begin{center}
{\color{darkaccent}\Large Алгоритм: как решить пропорцию\par}
\vspace{16mm}
\begin{tikzpicture}[
  >=Latex, line width=.65pt,
  start/.style={draw=darkaccent,rounded corners=7pt,minimum width=6.4cm,minimum height=1cm,align=center,fill=white},
  processbox/.style={draw=darkaccent,rounded corners=2pt,text width=9cm,minimum height=1.03cm,align=center,inner sep=7pt,fill=white},
  decision/.style={draw=darkaccent,diamond,aspect=3.6,text width=5.5cm,align=center,inner sep=1pt,fill=white},
  warning/.style={draw=linegray,rounded corners=2pt,text width=3.1cm,align=center,font=\small,inner sep=5pt},
  flow/.style={->,darkaccent,thick}
]
\node[start] (start) {Дана запись $\displaystyle\frac ab=\frac cd$};
\node[decision,below=10mm of start] (defined) {Знаменатели $b$ и $d$ не равны нулю?};
\node[warning,right=17mm of defined] (stop) {Нет: исходная дробная запись не определена};
\node[processbox,below=11mm of defined] (cross) {Перемножьте крайние и средние члены: $ad=bc$};
\node[processbox,below=7mm of cross] (brackets) {Сохраните скобки, если член пропорции --- сумма или произведение};
\node[processbox,below=7mm of brackets] (solve) {Решите полученное уравнение относительно неизвестной};
\node[processbox,below=7mm of solve] (verify) {Подставьте найденное значение в исходную пропорцию};
\node[start,below=8mm of verify] (finish) {Ответ};
\draw[flow] (start) -- (defined);
\draw[flow] (defined) -- node[right,font=\small] {да} (cross);
\draw[flow] (defined.east) -- node[above,font=\small] {нет} (stop.west);
\draw[flow] (cross) -- (brackets);
\draw[flow] (brackets) -- (solve);
\draw[flow] (solve) -- (verify);
\draw[flow] (verify) -- (finish);
\end{tikzpicture}
\end{center}
\clearpage

\section{Как выразить любую букву из пропорции}
Выразить переменную --- значит оставить её одну в левой части равенства,
а выражение через остальные величины записать справа.

Пусть $\dfrac ab=\dfrac cd$ и $b\ne0$, $d\ne0$.
По основному свойству пропорции $ad=bc$.

\begin{center}
\renewcommand{\arraystretch}{1.55}
\begin{tabular}{@{}ll@{}}
\toprule
\textcolor{darkaccent}{Что выразить} & \textcolor{darkaccent}{Результат}\\
\midrule
$a$ & $\displaystyle a=\frac{bc}{d}$\\
$b$ & $\displaystyle b=\frac{ad}{c}$, если $c\ne0$\\
$c$ & $\displaystyle c=\frac{ad}{b}$\\
$d$ & $\displaystyle d=\frac{bc}{a}$, если $a\ne0$\\
\bottomrule
\end{tabular}
\end{center}
\textbf{Откуда берётся формула?} Например, из $ad=bc$ получаем
$a=bc:d$, разделив обе части на $d$.

\textbf{Ограничения.} Делить на нуль нельзя. При выражении $b$ и $d$
нужно проверять дополнительные условия $c\ne0$ и $a\ne0$ соответственно.

\section{Выражение переменной из формулы по физике}
Изученный приём полезен и тогда, когда уравнение записано через буквы.

\example{Выразите $m_1$ из формулы}
\[
F=\frac{k m_1m_2}{r^2}.
\]
\begin{align*}
F r^2 &= k m_1m_2 &&\text{умножили обе части на $r^2$},\\
m_1(km_2)&=Fr^2 &&\text{собрали множители при $m_1$},\\
m_1&=\boxed{\frac{Fr^2}{km_2}} &&\text{разделили на $km_2$}.
\end{align*}
Предполагается $r\ne0$ и $km_2\ne0$.

\keyidea{Сложная запись с буквами подчиняется тем же правилам,
что и обычное уравнение с числами.}

\checkmarkline{Умею выразить любую из букв в пропорции при допустимых значениях.}
\checkmarkline{Умею убрать знаменатель и выделить искомый множитель.}

\section{Частые ошибки и самопроверка}
\begin{itemize}
\item \textbf{Перепутан порядок:} отношение $A$ к $B$ записывается как $A:B$.
\item \textbf{Не согласованы единицы:} $99$ г и $4{,}4$ кг нельзя делить как числа $99$ и $4{,}4$ без перевода.
\item \textbf{Неверно поняты доли:} в отношении $3:4$ доля первой группы от общего числа равна $\frac{3}{7}$, а не $\frac34$.
\item \textbf{Пропущены скобки:} $(a+0{,}03)$ следует сохранять как единый член пропорции.
\item \textbf{Забыто ограничение:} знаменатель дроби не может быть равен нулю.
\item \textbf{Не выполнена проверка:} подставьте ответ в исходное равенство.
\end{itemize}
\textcolor{darkaccent}{\textbf{Итоговая памятка}}\par
\begin{tabularx}{\linewidth}{@{}X X@{}}
\toprule
Сигнал в условии & Что сделать\\
\midrule
«Какую часть?» & Часть разделить на целое\\
«Во сколько раз?» & Первую величину разделить на вторую\\
«Разделите в отношении $p:q$» & Сумма долей $p+q$, затем одна доля\\
«Верна ли пропорция?» & Сравнить произведения крайних и средних\\
«Найдите неизвестный член» & Перемножить крест-накрест и решить уравнение\\
«Выразите переменную» & Убрать знаменатели, выделить нужную букву\\
\bottomrule
\end{tabularx}

\clearpage
\section*{Тренировочный блок}
\addcontentsline{toc}{section}{Тренировочный блок}
\textcolor{darkaccent}{\textbf{Путь А --- 10 заданий.}} Задания расположены по нарастанию сложности.\par
\textbf{Базовый уровень}
\begin{enumerate}[itemsep=5pt]
\item Из десяти задач решено семь. Какую часть составляют решённые задачи? Запишите ответ обыкновенной дробью и в процентах.
\item Из $450\ \text{м}^2$ площадки расчистили $270\ \text{м}^2$. Найдите:
  \begin{enumerate}[label=\alph*),nosep,topsep=2pt]
  \item какую часть площадки расчистили;
  \item во сколько раз вся площадка больше расчищенной части.
  \end{enumerate}
\item Найдите отношение $72$ см к $1{,}2$ м. Сократите полученную дробь.
\item Разделите $96$ конфет на две группы в отношении $5:3$. Сколько конфет будет в каждой группе?
\item Разделите число $72$ на три части в отношении $1:2:3$.
\end{enumerate}
\medskip\textbf{Уровень выше среднего}
\begin{enumerate}[start=6,itemsep=6pt]
\item Проверьте, верна ли пропорция $\displaystyle\frac58=\frac{15}{24}$. Обоснуйте ответ основным свойством пропорции.
\item Решите пропорцию $\displaystyle\frac{t}{9}=\frac8{12}$.
\item Решите пропорцию $\displaystyle\frac{7{,}2}{4}=\frac{3x}{5}$.
\end{enumerate}
\medskip\textbf{Повышенный уровень}
\begin{enumerate}[start=9,itemsep=6pt]
\item Решите пропорцию $\displaystyle\frac{1{,}8}{4{,}5}=\frac{a+0{,}06}{0{,}3}$.
\item Из формулы $\displaystyle F=\frac{k m_1m_2}{r^2}$ выразите $m_2$. Укажите, какие величины при этом должны быть ненулевыми.
\end{enumerate}
\clearpage

\section*{Ответы к тренировочному блоку}
\vspace{4mm}
\renewcommand{\arraystretch}{1.55}
\begin{tabularx}{\linewidth}{@{}p{16mm}X@{}}
\toprule
\textcolor{darkaccent}{№} & \textcolor{darkaccent}{Ответ}\\
\midrule
1 & $\frac7{10}=70\%$\\
2 & а) $\frac35$; б) $\frac53$\\
3 & $\frac35$\\
4 & $60$ и $36$ конфет\\
5 & $12$, $24$, $36$\\
6 & Верна: $5\cdot24=8\cdot15=120$\\
7 & $t=6$\\
8 & $x=3$\\
9 & $a=0{,}06$\\
10 & $\displaystyle m_2=\frac{Fr^2}{km_1}$; $r\ne0$, $km_1\ne0$\\
\bottomrule
\end{tabularx}

\vfill
\begin{center}
{\color{darkaccent}\small Лёвин Артём Александрович --- эксклюзивно для Нади Клименко}
\end{center}
\end{document}
